% Generated from tex/chapters/ch04/07_実務上の考慮.tex for the SWEBOK structure tree
\subsubsection{構築における再利用}

構築における再利用には、再利用のための構築と、再利用による構築がある。再利用のための構築では、将来別の場所で使えるように、部品を作る。可変性を設定で変えられるようにする、使い方を文書化する、単体テストを整備する、公開方法を整えるといった作業が必要である。

再利用による構築では、既存の部品を使って新しいソフトウェアを作る。ライブラリ、フレームワーク、オープンソース部品、商用既製品、クラウドサービスが代表例である。クラウドサービスを利用する場合、認証、メッセージング、ストレージなどを外部のサービスへ委任できる。このような形は、バックエンドサービスとしての利用とも呼ばれる。

再利用では、部品を選び、評価し、統合し、テストし、利用情報を記録する必要がある。ライセンス、保守状況、脆弱性、利用条件を確認しない再利用は、後のリスクになる。
